% --- ARCHIVO: Capitulos/introduccion.tex ---

\chapter{Sistema Lanza - Keitros}

\begin{quote}
\textit{''En este planeta llamado Keitros, extraños y desconocidos dioses poblaron estas tierras, y crearon seres de diferentes formas, cada dios se enfoco un grupo de ellos, y sin saber como o porque, cultivo a estos seres. Los Dioses les dieron poder e intelecto, sentir el poder. Y cuando estas creaturas se encontraron por primera vez en el mundo, ''La Era Del Martillo'' comenzó.''}
\end{quote}


\section{¿Qué es La Era del Martillo?}
\textbf{La Era del Martillo} es un juego de escaramuzas heroicas diseñado para jugadores que buscan profundidad táctica y consecuencias reales en cada decisión. A diferencia de otros juegos de guerra, aquí cada unidad es un activo invaluable. El sistema no se trata solo de mover miniaturas, sino de gestionar el cansancio, la posición y, sobre todo, la energía del entorno.

\begin{reglaespecial}{La Filosofía del Juego}
Este sistema se basa en tres pilares fundamentales:
\begin{itemize}
    \item \textbf{Letalidad Detallada:} Gracias al sistema de localización de daños, un solo golpe bien colocado puede cambiar el curso de la batalla.
    \item \textbf{Gestión de Recursos:} Los Puntos de Acción (PA) y Energía (PE) dictan el ritmo. Un guerrero agotado es un guerrero muerto.
    \item \textbf{Zonas Secas:} El uso de la magia tiene un costo físico en el campo de batalla, alterando el mapa de forma permanente.
\end{itemize}
\end{reglaespecial}

\section{Lo que Necesitas para Jugar}
Para comenzar tus batallas en la Era del Martillo, necesitarás lo siguiente:

\begin{description}
    \item[Miniaturas:] Que representen a tus héroes y escuadras.
    \item[Dados:] Un dado de doce caras (\textbf{1d12}) para pruebas de éxito, un dado de cien (\textbf{1d100}) para localización y un dado de seis (\textbf{1d6}) para el daño.
    \item[Superficie de Juego:] Un área de 90cm x 90cm con escenografía que represente ruinas, bosques o zonas de extracción de energía.
    \item[Cinta Métrica:] Las distancias se miden en centímetros (cm).
\end{description}

\section{El Concepto de Escaramuza}
A diferencia de las grandes batallas de bloques de infantería, este manual se enfoca en el **Combate Heroico**. Cada miniatura actúa de forma independiente. Una unidad puede ser un poderoso líder Evolgen o un pequeño grupo de apoyo Biskrorum, pero cada una tiene acceso a acciones, reacciones y habilidades únicas que consumen sus propios PA y PE.

\vfill
\begin{center}
    \textit{Prepárate, comandante. La Gema brilla, pero el terreno se seca.}
\end{center}